\documentclass[10pt,a4paper]{amsart}
\usepackage[margin=19mm]{geometry}
\usepackage{amsmath,amssymb,mathtools}
\usepackage[scr=rsfs,cal=euler]{mathalfa}
\usepackage{xfrac}
\usepackage[unicode,hidelinks,bookmarks=false]{hyperref}
\newcommand{\ie}{i.e.\ }
\newcommand{\cf}{cf.\ }
\newcommand{\resp}{respectively\ }
\newcommand{\Subsection}{\subsection}
\providecommand{\nobreakdash}{\nobreak\hbox{-}}
\setlength{\emergencystretch}{2em}
\multlinegap0pt
\usepackage{bm}
\usepackage{bbm}
\usepackage{dsfont}
\usepackage{xspace}
\usepackage{cancel}
\usepackage{mathtools}
\usepackage{amsthm}
\makeatletter
\@ifundefined{theorem}{\newtheorem{theorem}{Theorem}}{}
\@ifundefined{corollary}{\newtheorem{corollary}[theorem]{Corollary}}{}
\@ifundefined{lemma}{\newtheorem{lemma}[theorem]{Lemma}}{}
\makeatother
\def\CC{\mathscr{C}}
\def\E{\mathbb{E}}
\def\HH{\mathscr{H}}
\def\LL{\mathscr{L}}
\def\RR{\mathscr{R}}
\def\Y{\mathbb{Y}}
\def\ZZ{\mathscr Z}
\def\dom{{\sf D}}
\def\ran{{\sf R}}
\title[A combinatorial approach to phase transitions in random grap]{A combinatorial approach to phase transitions in random graph isomorphism problems}
\author{Dimitris Diamantidis, Takis Konstantopoulos, Linglong Yuan}
\date{Source: arXiv:2410.00214v2 (CC BY 4.0)}
\begin{document}
\maketitle

\noindent\emph{Source and licence.} A combinatorial approach to phase transitions in random graph isomorphism problems. Version arXiv:2410.00214v2 (CC BY 4.0), distributed under the Creative Commons Attribution 4.0 International licence, as recorded in the arXiv licence metadata for this exact version and shown on the abstract page https://arxiv.org/abs/2410.00214v2. Full rights notice: licenses/arxiv-2410.00214-CC-BY-4.0.txt.\par\smallskip

\noindent\textit{Editorial scope.} Counted unit: the Lemma whose source label is \texttt{shuplem} in the article (source0.tex lines 2503--2533, source label \texttt{shuplem}), delivered together with the article's own complete original proof of it (source0.tex lines 2535--2660, one \texttt{proof} environment of 126 lines). No mathematics is written, repaired or re-derived: statement and proof are the article's own text line for line. Required context retained uncounted: thespis (lines 772--776); ccut (lines 2011--2018); N2tris (lines 2101--2104); right (lines 2185--2204); por1 (lines 2232--2254); por3 (lines 2233--2253); por4 (lines 2233--2253); lambduplem (lines 2263--2269); lambdupeq (lines 2265--2268). Every definition, display and label of those lines is retained unchanged. Omitted from this package and not redistributed: 1--771, 777--2010, 2019--2100, 2105--2184, 2205--2231, 2254--2262, 2269--2502, 2534--2534, 2661--3644. Nothing inside a retained excerpt is omitted or altered: every retained body is delivered exactly as the article's lines are written, so no omission receipt of any kind accompanies this package. Citations: the retained excerpts carry no citation command at all, so no bibliography entry of the article is transcribed and no reference is redistributed. Reference adaptation: none was needed and none was made; every reference inside the delivered unit resolves inside it, so no unresolved reference or citation is produced. Printed slips kept literal: no printed word, symbol, hypothesis or number of the counted statement or of its proof was altered. Layout and class: the article's own class is \texttt{amsart} with packages including amssymb, graphicx, lipsum, geometry, xcolor, hyperref, caption, mathrsfs, enumerate, bold-extra, cancel, ulem, datetime2, cases; the wrapper is the lane's pinned \texttt{amsart} editorial wrapper at 10pt on A4, which restates the theorem-like environments the retained text needs and adds the article's own macro definitions that the retained text needs (\texttt{\textbackslash CC}, \texttt{\textbackslash E}, \texttt{\textbackslash HH}, \texttt{\textbackslash LL}, \texttt{\textbackslash RR}, \texttt{\textbackslash Y}, \texttt{\textbackslash ZZ}, \texttt{\textbackslash dom}, \texttt{\textbackslash ran}), with extra wrapper packages (bm, bbm, dsfont, xspace, cancel, mathtools). The delivered numbering is the unit's own. Assets: no retained span contains a figure environment, a graphics-inclusion command, a PDF-page inclusion, a table or a TikZ picture; no image file is copied into this package, the package declares no asset dependency and no image file is redistributed with it. Licence: version arXiv:2410.00214v2 of this article is distributed under the Creative Commons Attribution 4.0 International licence, as recorded in the arXiv licence metadata for this exact version and shown on the abstract page https://arxiv.org/abs/2410.00214v2. A full rights notice is delivered at licenses/arxiv-2410.00214-CC-BY-4.0.txt. Characters: every retained excerpt is pure ASCII, so the delivered engine, XeLaTeX with its default Latin Modern text font, reports no missing character.
\begin{equation}
\label{thespis}
\Y:=\left\{(p,q) \in (0,1)\times (0,1):\, \max\{\tau_{1,2}(p,q), \tau_{2,1}(p,q)\right\} 
< \tau(p,q)^{3/2}\},
\end{equation}

\begin{lemma}
\label{ccut}
(i) The set $\CC_{j,k}(f,g)$ is empty unless $j=k-1$ or $j=k$ or $j=k+1$.
\\
(ii) If $C \in \CC_{j,k}(f,g)$ with $|k-j|=1$ then $C$ is a path.
\\
(iii) If $C \in \CC_{j,j}(f,g)$ then $C$ is a path or a cycle.
\end{lemma}

\begin{equation}
\label{N2tris} 
\E J_f J_g = \prod_{j,k \ge 1} \tau_{j,k}^{|\CC_{j,k}(f,g)|}.
\end{equation}

\begin{align}
|\LL| &= 
\sum_{j\ge 1} (j |\CC_{j,j}| + j |\CC_{j,j+1}| + (j+1) |\CC_{j+1,j}| )
= 2 \binom{m}{2} - \binom{|Z|}{2},
\label{left}
\\
|\RR|&=\sum_{j\ge 1} (j |\CC_{j,j}| + j |\CC_{j+1,j}| + (j+1) |\CC_{j,j+1}| )
= 2 \binom{m}{2} - \binom{|\ZZ|}{2}.
\label{right}
\\
|\LL_1|&=\sum_{j\ge 1} (|\CC_{j,j}^*| + 2 |\CC_{j+1,j}|)
= |\ZZ| +  2 \binom{m}{2} - 2 \binom{|Z|}{2},
%\qquad \text{ no.\ of left VERTICES of DEG 1}
\label{left1}
\\
|\RR_1|&=\sum_{j\ge 1} (|\CC_{j,j}^*| + 2 |\CC_{j,j+1}|)
= |\ZZ| +  2 \binom{m}{2} - 2 \binom{\ZZ}{2}.
%\qquad \text{ no.\ of right VERTICES of DEG 1}
\label{right1}
\end{align}

\begin{corollary}
\begin{align}
&\sum_{j \ge 1} \left\{ |\CC_{j,j+1}| - |\CC_{j+1,j}|\right\} 
= \binom{d}{2}-\binom{r}{2}
\label{por1}
\\
&\sum_{j \ge 1} \left\{(j-1)|\CC_{j,j}^*| + j |\CC_{j,j}^o| 
+ (j-1)|\CC_{j,j+1}| + j |\CC_{j+1,j}| \right\} 
= \binom{r}{2} - |\ZZ|
\label{por2}
\\
&\sum_{j \ge 2} |\CC_{j, j}^o| \le \frac{1}{2} \left( \binom{r}{2} - |\ZZ| \right)
\label{por3}
\\
&
%\sum_{j \ge 1} ((j-1) (|\CC_{j, j}| + |\CC_{j, j+1}|) + j |\CC_{j+1, j}|)
%\ge \frac{1}{2} \left( \binom{r}{2} - |\ZZ| \right)
\sum_{j \ge 1} \left\{
(j-1) |\CC_{j, j}| + (j-1)|\CC_{j, j+1}| + j |\CC_{j+1, j}| \right\}
\ge \frac{1}{2} \left( \binom{r}{2} - |\ZZ| \right)
\label{por4}
\end{align}
\end{corollary}

\begin{lemma} 
\label{lambduplem}
\begin{equation}
\label{lambdupeq}
\binom{|Z|}{2}  \le |\ZZ| \le \binom{|Z|}{2} + \frac{1}{2} (r - |Z|).
\end{equation}
\end{lemma}

\begin{lemma}
\label{shuplem}
Assume that condition in \eqref{thespis} holds:
\[
\tau^{3/2} > \max(\tau_{1,2}, \tau_{2,1}).
\]
Set
\begin{equation}
\omega := \frac{\max\{pq, (1-p)(1-q)\}}{pq + (1-p)(1-q)}, 
\quad
\beta := \sqrt{\max\left\{\omega, \frac{\tau_{1, 2}\tau_{2, 1}}{\tau^3}\right\}},
\quad
\gamma := \lambda \log (\tau/\tau_{1,2})
\label{gammadef}
\end{equation} 
Then
\[
0 \le \beta < 1, \qquad \frac{1}{2} < \gamma \le 1,
\]
and, if $r \le d$ then,
for all $(f,g) \in \HH_{d,r,\ell}$, $0 \le \ell \le r$,
\begin{equation} 
\label{shupeq}
\frac{\E J_f J_g}{\E J_f \E J_g}
%\le \left( \frac{1}{\tau} \right)^{\binom{d}{2}}
%\left( \frac{\tau_{1,2}}{\tau} \right)^{\binom{d}{2} - \binom{r}{2}} 
%\beta^{\frac{1}{2}(r - \ell)(r-2)}
\le  \left( \frac{1}{\tau} \right)^{(1-\gamma)\binom{d}{2} + \gamma\binom{r}{2}}
\beta^{\frac{1}{2}(r - \ell)(r-2)}.
\end{equation}
\end{lemma}

\begin{proof}
Since $\omega<1$ and
$\tau_{1,2} \tau_{2,1} \le \max(\tau_{1,2}, \tau_{2,1})^2 < \tau^3$. 
we have that $\beta<1$.
Next, by elementary algebra, $\tau^2 \le \tau_{1,2}$. Hence $\tau/\tau_{1/2}  \le 1/\tau$,
and, taking logarithms, $\log(\tau/\tau_{1/2}) \le \log(1/\tau)=1/\lambda$.
This gives $\gamma \le 1$.
Since the condition in \eqref{thespis} holds, $\tau^{3/2} >\tau_{1,2}$. 
Hence $\tau/\tau_{1,2} >1/\tau^{1/2}$.
Taking logarithms we obtain $\log(\tau/\tau_{1,2}) > \frac12 \log(1/\tau)
=\frac{1}{2\lambda}$.
This gives $\gamma >1/2$.
Assume $(f,g) \in \HH_{d,r,\ell}$, which means that
$\dom f \cap \dom g$ has size $d$, $\ran f \cap \ran g$ has size $r$ 
and $Z(f,g)$ has size $\ell$.
We now use the expression 
$\E J_f J_g=\prod_{j,k \ge 1} \tau_{j,k}^{|\CC_{j,k}(f,g)|}$ obtained in
Proposition \ref{N2tris}.
Write $\bar p = 1-p$, $\bar q = 1-q$.
We write 
\begin{align*}
\tau_{j,k} 
&= (pq)^{j-1} p q^{k-j+1} 
+ (\bar p \bar q)^{j-1}  \bar p  \bar q^{k-j+1} 
\\
&\le \max\left\{(pq)^{j-1},\, (\bar p\bar q)^{j-1}\right\}
\, (p q^{k-j+1} + \bar p \bar q^{k-j+1}) 
\\
& = (\omega \tau)^{j-1} \, (p q^{k-j+1} + \bar p \bar q^{k-j+1}).
\end{align*}
We obtain a second inequality by interchanging $j$ and $k$.
We therefore have
\[
\tau_{j,k} \le \begin{cases}
(\omega \tau)^{j-1} \tau_{1,k-j+1} & \text{ if } j \le k
\\
(\omega \tau)^{k-1} \tau_{j-k+1,1} & \text{ if } k \le j
\end{cases}.
\]
Using these inequalities in the expression for $\E J_f J_g$ we have
\begin{align*}
\E J_f J_g \le
\prod_{\substack{j,k \ge 1\\ j \le k}} 
\left[(\omega \tau)^{j-1} \tau_{1,k-j+1}\right]^{|\CC_{j,k}(f,g)|}
\prod_{\substack{j,k \ge 1\\ k<j}} 
\left[(\omega \tau)^{k-1} \tau_{j-k+1,1}\right]^{|\CC_{j,k}(f,g)|}.
\end{align*}
By Lemma \ref{ccut}, only the terms of the form $(j,j+1)$ or $(j,j)$ 
in the first product, and only the terms of the form $(k,k+1)$ in the second product,
survive. Making the change of variable $(j,k) \to (k,j)$ in the second
product and grouping terms together, we obtain
\begin{equation}
\label{once}
\E J_f J_g \le
\omega^{\mathsf A} \, \tau^{\mathsf B}
\, \tau_{1,2}^{\mathsf D_{1,2}}
\, \tau_{2,1}^{\mathsf D_{2,1}},
\end{equation}
where
\begin{align*}
\mathsf A &= \sum_{j \ge 1} 
\big\{(j-1) |\CC_{j,j+1}| + (j-1) |\CC_{j+1,j}| + (j-1) |\CC_{j,j}| \big\}
\\
\mathsf B &= \sum_{j \ge 1} 
\big\{(j-1) |\CC_{j,j+1}| + (j-1) |\CC_{j+1,j}| + j |\CC_{j,j}| \big\}
\\ & \mathsf D_{1,2} = \sum_{j \ge 1} |\CC_{j,j+1}|, \quad
\mathsf D_{2,1} = \sum_{j \ge 1} |\CC_{j+1,j}|
\end{align*}
Using these symbols, \eqref{right} reads
\[
\mathsf B + \mathsf D_{1,2} + 2 \mathsf D_{2,1} = 2 \binom{m}{2}-\binom{d}{2},
\]
and so the right-hand side of \eqref{once} equals
\begin{align}
\omega^{\mathsf A} \tau^{2 \binom{m}{2}-\binom{d}{2}}
\left(\frac{\tau_{1,2}}{\tau}\right)^{\mathsf D_{1,2}} 
\left(\frac{\tau_{2,1}}{\tau^2}\right)^{\mathsf D_{2,1}} 
=
\omega^{\mathsf A} \tau^{2 \binom{m}{2}-\binom{d}{2}}
\left(\frac{\tau_{1,2}}{\tau}\right)^{\mathsf D_{1,2}-\mathsf D_{2,1}}
\left(\frac{\tau_{1,2}\tau_{2,1}}{\tau^3}\right)^{\mathsf D_{2,1}} 
\nonumber
\\
\le \beta^{2 \mathsf A+2 \mathsf D_{1,2}} \tau^{2 \binom{m}{2}-\binom{d}{2}}
\left(\frac{\tau_{1,2}}{\tau}\right)^{\mathsf D_{1,2}-\mathsf D_{2,1}} .
\label{twice}
\end{align}
Rewrite \eqref{por1} and \eqref{por4} as
\[
\mathsf D_{1,2}-\mathsf D_{2,1} = \binom{d}{2}-\binom{r}{2},
\quad
2 \mathsf A+2 \mathsf D_{1,2} \ge \binom{r}{2}-|\ZZ|.
\]
Note that both terms are nonnegative, the first due to the assumption
$r \le d$, and the second due to \eqref{por3}.
We then see that the right-hand side of 
\eqref{twice} is majorized by the right-hand side
of:
\begin{equation}
\label{thrice}
\E J_f J_g 
\le 
\beta^{\binom{r}{2}-|\ZZ|} \tau^{2\binom{m}{2}-\binom{d}{2}}
%\left(\frac{1}{\tau}\right)^{\binom{d}{2}}
\left(\frac{\tau_{1,2}}{\tau}\right)^{\binom{d}{2}-\binom{r}{2}}.
\end{equation}
To conclude the proof, we use the second inequality 
in \eqref{lambdupeq} of Lemma \eqref{lambduplem} to obtain
\[
\binom{r}{2}-|\ZZ| \ge \binom{|Z|}{2}+\frac12 (r-|Z|)
= \binom{r}{2}- \binom{\ell}{2}+\frac12 (r-\ell)
\ge \frac12 (r-\ell)(r-2),
\]
Replacing the exponent of $\beta$ in \eqref{thrice} by the latter
quantity and
dividing both sides by $(\E J_f)(\E J_g) = \tau^{2\binom{m}{2}}$,
we obtain at \eqref{shupeq}.
\[
\frac{\E J_f J_g}{\E J_f \E J_g}
\le \left( \frac{1}{\tau} \right)^{\binom{d}{2}}
\left( \frac{\tau_{1,2}}{\tau} \right)^{\binom{d}{2} - \binom{r}{2}} 
\beta^{\frac{1}{2}(r - \ell)(r-2)}
=  \left( \frac{1}{\tau} \right)^{(1-\gamma)\binom{d}{2} + \gamma\binom{r}{2}},
\]
by the definition of $\gamma$.
\end{proof}

\end{document}
